% ============================================================
% Integrante: Ernesto A. Soler Choong
% Fragmento puro: se incluye desde preguntas/p5-kitty-problem/solucion.tex
% mediante % ============================================================
% Integrante: Ernesto A. Soler Choong
% Fragmento puro: se incluye desde preguntas/p5-kitty-problem/solucion.tex
% mediante % ============================================================
% Integrante: Ernesto A. Soler Choong
% Fragmento puro: se incluye desde preguntas/p5-kitty-problem/solucion.tex
% mediante % ============================================================
% Integrante: Ernesto A. Soler Choong
% Fragmento puro: se incluye desde preguntas/p5-kitty-problem/solucion.tex
% mediante \input{investigacion/p5-kitty-ocp/notas}.
% Título: se omite «en Richards & Ford» del título oficial porque el Kitty
% Fuente de trabajo: notas.md (misma carpeta).
% Citas: páginas impresas; Richards & Ford (2020) no tiene paginación y se
% cita por capítulo/sección. Todas las claves están en core.bib.
% ============================================================

% Grupo local: \emergencystretch evita líneas desbordadas sin afectar a otros capítulos.
\begingroup\emergencystretch=3em \hbadness=2000
\section*{Notas de investigación: Kitty Problem, OCP y Plugin Architecture}

\section{Síntesis}

El \emph{Kitty Problem} es un ejemplo de \textcite[242--243]{martin2017}, no de Richards y
Ford. Parte de un agregador de taxis, presentado antes en el capítulo sobre el Principio
de Sustitución de Liskov \parencite[80]{martin2017}, construido como microservicios
divididos por función: \texttt{TaxiUI}, \texttt{TaxiFinder}, \texttt{TaxiSelector} y
\texttt{TaxiDispatcher} \parencite[242]{martin2017}. Cuando se pide un servicio de entrega
de gatitos, con restricciones por alergias de conductores y clientes, hay que modificar
\emph{todos} los servicios y coordinar su despliegue \parencite[243]{martin2017}. Martin
usa el caso para desmontar dos supuestos beneficios de los servicios: el desacoplamiento
y el desarrollo y despliegue independientes \parencite[240--241]{martin2017}. Su
diagnóstico es que las descomposiciones funcionales son muy vulnerables a los aspectos
transversales \parencite[243]{martin2017}.

La salida que propone no es crear más servicios, sino aplicar SOLID dentro de ellos
\parencite[244]{martin2017}. La lógica genérica queda en clases base abstractas; lo
específico de los viajes se extrae a un componente \texttt{Rides} y la nueva funcionalidad
va a un componente \texttt{Kittens}. Ambos extienden esas abstracciones mediante
\emph{Template Method} o \emph{Strategy}, respetan la Regla de Dependencia y son creados
por factorías controladas por la interfaz de usuario. Solo cambia \texttt{TaxiUI}; el
resto se extiende añadiendo un componente que se carga en tiempo de ejecución
\parencite[244]{martin2017}, lo que es el Principio Abierto/Cerrado aplicado a
componentes \parencite[246]{martin2017}. La lección arquitectónica: las fronteras no
pasan \emph{entre} los servicios, sino \emph{a través} de ellos
\parencite[246--247]{martin2017}.

Completan el cuadro tres ideas. El CCP es el SRP llevado a componentes, y Martin lo
vincula expresamente con el \emph{Kitty Problem} \parencite[105--106]{martin2017}. La
arquitectura de \emph{plugins} aísla las reglas de negocio de los componentes opcionales
o intercambiables \parencite[170--173]{martin2017}. Y
\textcite[caps.~8 y 17]{richardsford2020} aportan la partición por dominio, el antipatrón
\emph{entity trap} y el bounded context como criterio de granularidad.

\section{Conceptos clave}

\begin{table}[htbp]
  \centering
  \footnotesize
  \renewcommand{\arraystretch}{1.2}
  \begin{tabular}{@{}p{0.21\linewidth}p{0.47\linewidth}p{0.24\linewidth}@{}}
    \toprule
    \textbf{Concepto} & \textbf{Definición} & \textbf{Fuente} \\
    \midrule
    \raggedright \emph{Kitty Problem}
      & \raggedright Un requisito transversal obliga a cambiar y desplegar de forma
        coordinada todos los servicios de un sistema descompuesto por función
      & \raggedright \textcite[242--243]{martin2017} \tabularnewline
    \raggedright Falacias de los servicios
      & \raggedright Siguen acoplados por los datos que comparten; en esa medida, su
        desarrollo y despliegue deben coordinarse
      & \raggedright \textcite[240--241]{martin2017} \tabularnewline
    \raggedright OCP
      & \raggedright Abierto a la extensión y cerrado a la modificación (Meyer, citado por
        Martin); en arquitectura, con SRP y DIP
      & \raggedright \textcite[70, 73--74]{martin2017} \tabularnewline
    \raggedright CCP y cierre estratégico
      & \raggedright Reunir lo que cambia por las mismas razones; cerrar frente a los
        cambios más habituales
      & \raggedright \textcite[105--106]{martin2017} \tabularnewline
    \raggedright Arquitectura de \emph{plugins}
      & \raggedright Las reglas de negocio no dependen de los componentes opcionales; las
        fronteras siguen los ejes de cambio
      & \raggedright \textcite[170--173]{martin2017} \tabularnewline
    \raggedright Regla de Dependencia y adaptadores
      & \raggedright Las dependencias apuntan hacia adentro; los adaptadores convierten
        datos entre el núcleo y lo externo
      & \raggedright \textcite[203, 205]{martin2017} \tabularnewline
    \raggedright Servicios basados en componentes
      & \raggedright Estructura interna SOLID; lo nuevo entra como componentes nuevos y las
        fronteras atraviesan los servicios
      & \raggedright \textcite[245--247]{martin2017} \tabularnewline
    \raggedright Patrón microkernel
      & \raggedright Núcleo mínimo más \emph{plugins} independientes, localizados mediante
        un registro; adaptador para contratos propios
      & \raggedright \textcite[21--23]{richards2015} \tabularnewline
    \raggedright Partición por dominio y \emph{entity trap}
      & \raggedright Con partición técnica, un flujo se reparte entre capas; un componente
        por entidad confunde datos con flujos
      & \raggedright \textcite[cap.~8]{richardsford2020} \tabularnewline
    \raggedright Bounded context y granularidad
      & \raggedright Cada servicio modela un dominio o flujo; fronteras por propósito,
        transacciones y coreografía
      & \raggedright \textcite[caps.~7 y 17]{richardsford2020} \tabularnewline
    \bottomrule
  \end{tabular}
  \caption{Conceptos clave de la investigación INV-P5.}
  \label{tab:inv-p5-conceptos}
\end{table}

\section{Puntos de conexión con las preguntas}

\begin{itemize}
  \item \textbf{P5, por qué falló dividir por verbos o entidades:} la vulnerabilidad de la
        descomposición funcional, las falacias de los servicios, la partición técnica y el
        \emph{entity trap}, y la violación del CCP.
  \item \textbf{P5, diseño con DDD y Onion/Clean:} fronteras sobre los ejes de cambio,
        núcleo protegido por la Regla de Dependencia y bounded context como criterio de
        granularidad.
  \item \textbf{P5, \texttt{Kittens} y \texttt{Rides} como \emph{plugins}:} la solución
        de Martin, el OCP aplicado a componentes, el patrón microkernel y los adaptadores.
  \item \textbf{P3 y P4:} el bounded context y la granularidad apoyan la división en
        contextos de P3; las clases base abstractas y las factorías ilustran el DIP de P4.
\end{itemize}


.
% Título: se omite «en Richards & Ford» del título oficial porque el Kitty
% Fuente de trabajo: notas.md (misma carpeta).
% Citas: páginas impresas; Richards & Ford (2020) no tiene paginación y se
% cita por capítulo/sección. Todas las claves están en core.bib.
% ============================================================

% Grupo local: \emergencystretch evita líneas desbordadas sin afectar a otros capítulos.
\begingroup\emergencystretch=3em \hbadness=2000
\section*{Notas de investigación: Kitty Problem, OCP y Plugin Architecture}

\section{Síntesis}

El \emph{Kitty Problem} es un ejemplo de \textcite[242--243]{martin2017}, no de Richards y
Ford. Parte de un agregador de taxis, presentado antes en el capítulo sobre el Principio
de Sustitución de Liskov \parencite[80]{martin2017}, construido como microservicios
divididos por función: \texttt{TaxiUI}, \texttt{TaxiFinder}, \texttt{TaxiSelector} y
\texttt{TaxiDispatcher} \parencite[242]{martin2017}. Cuando se pide un servicio de entrega
de gatitos, con restricciones por alergias de conductores y clientes, hay que modificar
\emph{todos} los servicios y coordinar su despliegue \parencite[243]{martin2017}. Martin
usa el caso para desmontar dos supuestos beneficios de los servicios: el desacoplamiento
y el desarrollo y despliegue independientes \parencite[240--241]{martin2017}. Su
diagnóstico es que las descomposiciones funcionales son muy vulnerables a los aspectos
transversales \parencite[243]{martin2017}.

La salida que propone no es crear más servicios, sino aplicar SOLID dentro de ellos
\parencite[244]{martin2017}. La lógica genérica queda en clases base abstractas; lo
específico de los viajes se extrae a un componente \texttt{Rides} y la nueva funcionalidad
va a un componente \texttt{Kittens}. Ambos extienden esas abstracciones mediante
\emph{Template Method} o \emph{Strategy}, respetan la Regla de Dependencia y son creados
por factorías controladas por la interfaz de usuario. Solo cambia \texttt{TaxiUI}; el
resto se extiende añadiendo un componente que se carga en tiempo de ejecución
\parencite[244]{martin2017}, lo que es el Principio Abierto/Cerrado aplicado a
componentes \parencite[246]{martin2017}. La lección arquitectónica: las fronteras no
pasan \emph{entre} los servicios, sino \emph{a través} de ellos
\parencite[246--247]{martin2017}.

Completan el cuadro tres ideas. El CCP es el SRP llevado a componentes, y Martin lo
vincula expresamente con el \emph{Kitty Problem} \parencite[105--106]{martin2017}. La
arquitectura de \emph{plugins} aísla las reglas de negocio de los componentes opcionales
o intercambiables \parencite[170--173]{martin2017}. Y
\textcite[caps.~8 y 17]{richardsford2020} aportan la partición por dominio, el antipatrón
\emph{entity trap} y el bounded context como criterio de granularidad.

\section{Conceptos clave}

\begin{table}[htbp]
  \centering
  \footnotesize
  \renewcommand{\arraystretch}{1.2}
  \begin{tabular}{@{}p{0.21\linewidth}p{0.47\linewidth}p{0.24\linewidth}@{}}
    \toprule
    \textbf{Concepto} & \textbf{Definición} & \textbf{Fuente} \\
    \midrule
    \raggedright \emph{Kitty Problem}
      & \raggedright Un requisito transversal obliga a cambiar y desplegar de forma
        coordinada todos los servicios de un sistema descompuesto por función
      & \raggedright \textcite[242--243]{martin2017} \tabularnewline
    \raggedright Falacias de los servicios
      & \raggedright Siguen acoplados por los datos que comparten; en esa medida, su
        desarrollo y despliegue deben coordinarse
      & \raggedright \textcite[240--241]{martin2017} \tabularnewline
    \raggedright OCP
      & \raggedright Abierto a la extensión y cerrado a la modificación (Meyer, citado por
        Martin); en arquitectura, con SRP y DIP
      & \raggedright \textcite[70, 73--74]{martin2017} \tabularnewline
    \raggedright CCP y cierre estratégico
      & \raggedright Reunir lo que cambia por las mismas razones; cerrar frente a los
        cambios más habituales
      & \raggedright \textcite[105--106]{martin2017} \tabularnewline
    \raggedright Arquitectura de \emph{plugins}
      & \raggedright Las reglas de negocio no dependen de los componentes opcionales; las
        fronteras siguen los ejes de cambio
      & \raggedright \textcite[170--173]{martin2017} \tabularnewline
    \raggedright Regla de Dependencia y adaptadores
      & \raggedright Las dependencias apuntan hacia adentro; los adaptadores convierten
        datos entre el núcleo y lo externo
      & \raggedright \textcite[203, 205]{martin2017} \tabularnewline
    \raggedright Servicios basados en componentes
      & \raggedright Estructura interna SOLID; lo nuevo entra como componentes nuevos y las
        fronteras atraviesan los servicios
      & \raggedright \textcite[245--247]{martin2017} \tabularnewline
    \raggedright Patrón microkernel
      & \raggedright Núcleo mínimo más \emph{plugins} independientes, localizados mediante
        un registro; adaptador para contratos propios
      & \raggedright \textcite[21--23]{richards2015} \tabularnewline
    \raggedright Partición por dominio y \emph{entity trap}
      & \raggedright Con partición técnica, un flujo se reparte entre capas; un componente
        por entidad confunde datos con flujos
      & \raggedright \textcite[cap.~8]{richardsford2020} \tabularnewline
    \raggedright Bounded context y granularidad
      & \raggedright Cada servicio modela un dominio o flujo; fronteras por propósito,
        transacciones y coreografía
      & \raggedright \textcite[caps.~7 y 17]{richardsford2020} \tabularnewline
    \bottomrule
  \end{tabular}
  \caption{Conceptos clave de la investigación INV-P5.}
  \label{tab:inv-p5-conceptos}
\end{table}

\section{Puntos de conexión con las preguntas}

\begin{itemize}
  \item \textbf{P5, por qué falló dividir por verbos o entidades:} la vulnerabilidad de la
        descomposición funcional, las falacias de los servicios, la partición técnica y el
        \emph{entity trap}, y la violación del CCP.
  \item \textbf{P5, diseño con DDD y Onion/Clean:} fronteras sobre los ejes de cambio,
        núcleo protegido por la Regla de Dependencia y bounded context como criterio de
        granularidad.
  \item \textbf{P5, \texttt{Kittens} y \texttt{Rides} como \emph{plugins}:} la solución
        de Martin, el OCP aplicado a componentes, el patrón microkernel y los adaptadores.
  \item \textbf{P3 y P4:} el bounded context y la granularidad apoyan la división en
        contextos de P3; las clases base abstractas y las factorías ilustran el DIP de P4.
\end{itemize}


.
% Título: se omite «en Richards & Ford» del título oficial porque el Kitty
% Fuente de trabajo: notas.md (misma carpeta).
% Citas: páginas impresas; Richards & Ford (2020) no tiene paginación y se
% cita por capítulo/sección. Todas las claves están en core.bib.
% ============================================================

% Grupo local: \emergencystretch evita líneas desbordadas sin afectar a otros capítulos.
\begingroup\emergencystretch=3em \hbadness=2000
\section*{Notas de investigación: Kitty Problem, OCP y Plugin Architecture}

\section{Síntesis}

El \emph{Kitty Problem} es un ejemplo de \textcite[242--243]{martin2017}, no de Richards y
Ford. Parte de un agregador de taxis, presentado antes en el capítulo sobre el Principio
de Sustitución de Liskov \parencite[80]{martin2017}, construido como microservicios
divididos por función: \texttt{TaxiUI}, \texttt{TaxiFinder}, \texttt{TaxiSelector} y
\texttt{TaxiDispatcher} \parencite[242]{martin2017}. Cuando se pide un servicio de entrega
de gatitos, con restricciones por alergias de conductores y clientes, hay que modificar
\emph{todos} los servicios y coordinar su despliegue \parencite[243]{martin2017}. Martin
usa el caso para desmontar dos supuestos beneficios de los servicios: el desacoplamiento
y el desarrollo y despliegue independientes \parencite[240--241]{martin2017}. Su
diagnóstico es que las descomposiciones funcionales son muy vulnerables a los aspectos
transversales \parencite[243]{martin2017}.

La salida que propone no es crear más servicios, sino aplicar SOLID dentro de ellos
\parencite[244]{martin2017}. La lógica genérica queda en clases base abstractas; lo
específico de los viajes se extrae a un componente \texttt{Rides} y la nueva funcionalidad
va a un componente \texttt{Kittens}. Ambos extienden esas abstracciones mediante
\emph{Template Method} o \emph{Strategy}, respetan la Regla de Dependencia y son creados
por factorías controladas por la interfaz de usuario. Solo cambia \texttt{TaxiUI}; el
resto se extiende añadiendo un componente que se carga en tiempo de ejecución
\parencite[244]{martin2017}, lo que es el Principio Abierto/Cerrado aplicado a
componentes \parencite[246]{martin2017}. La lección arquitectónica: las fronteras no
pasan \emph{entre} los servicios, sino \emph{a través} de ellos
\parencite[246--247]{martin2017}.

Completan el cuadro tres ideas. El CCP es el SRP llevado a componentes, y Martin lo
vincula expresamente con el \emph{Kitty Problem} \parencite[105--106]{martin2017}. La
arquitectura de \emph{plugins} aísla las reglas de negocio de los componentes opcionales
o intercambiables \parencite[170--173]{martin2017}. Y
\textcite[caps.~8 y 17]{richardsford2020} aportan la partición por dominio, el antipatrón
\emph{entity trap} y el bounded context como criterio de granularidad.

\section{Conceptos clave}

\begin{table}[htbp]
  \centering
  \footnotesize
  \renewcommand{\arraystretch}{1.2}
  \begin{tabular}{@{}p{0.21\linewidth}p{0.47\linewidth}p{0.24\linewidth}@{}}
    \toprule
    \textbf{Concepto} & \textbf{Definición} & \textbf{Fuente} \\
    \midrule
    \raggedright \emph{Kitty Problem}
      & \raggedright Un requisito transversal obliga a cambiar y desplegar de forma
        coordinada todos los servicios de un sistema descompuesto por función
      & \raggedright \textcite[242--243]{martin2017} \tabularnewline
    \raggedright Falacias de los servicios
      & \raggedright Siguen acoplados por los datos que comparten; en esa medida, su
        desarrollo y despliegue deben coordinarse
      & \raggedright \textcite[240--241]{martin2017} \tabularnewline
    \raggedright OCP
      & \raggedright Abierto a la extensión y cerrado a la modificación (Meyer, citado por
        Martin); en arquitectura, con SRP y DIP
      & \raggedright \textcite[70, 73--74]{martin2017} \tabularnewline
    \raggedright CCP y cierre estratégico
      & \raggedright Reunir lo que cambia por las mismas razones; cerrar frente a los
        cambios más habituales
      & \raggedright \textcite[105--106]{martin2017} \tabularnewline
    \raggedright Arquitectura de \emph{plugins}
      & \raggedright Las reglas de negocio no dependen de los componentes opcionales; las
        fronteras siguen los ejes de cambio
      & \raggedright \textcite[170--173]{martin2017} \tabularnewline
    \raggedright Regla de Dependencia y adaptadores
      & \raggedright Las dependencias apuntan hacia adentro; los adaptadores convierten
        datos entre el núcleo y lo externo
      & \raggedright \textcite[203, 205]{martin2017} \tabularnewline
    \raggedright Servicios basados en componentes
      & \raggedright Estructura interna SOLID; lo nuevo entra como componentes nuevos y las
        fronteras atraviesan los servicios
      & \raggedright \textcite[245--247]{martin2017} \tabularnewline
    \raggedright Patrón microkernel
      & \raggedright Núcleo mínimo más \emph{plugins} independientes, localizados mediante
        un registro; adaptador para contratos propios
      & \raggedright \textcite[21--23]{richards2015} \tabularnewline
    \raggedright Partición por dominio y \emph{entity trap}
      & \raggedright Con partición técnica, un flujo se reparte entre capas; un componente
        por entidad confunde datos con flujos
      & \raggedright \textcite[cap.~8]{richardsford2020} \tabularnewline
    \raggedright Bounded context y granularidad
      & \raggedright Cada servicio modela un dominio o flujo; fronteras por propósito,
        transacciones y coreografía
      & \raggedright \textcite[caps.~7 y 17]{richardsford2020} \tabularnewline
    \bottomrule
  \end{tabular}
  \caption{Conceptos clave de la investigación INV-P5.}
  \label{tab:inv-p5-conceptos}
\end{table}

\section{Puntos de conexión con las preguntas}

\begin{itemize}
  \item \textbf{P5, por qué falló dividir por verbos o entidades:} la vulnerabilidad de la
        descomposición funcional, las falacias de los servicios, la partición técnica y el
        \emph{entity trap}, y la violación del CCP.
  \item \textbf{P5, diseño con DDD y Onion/Clean:} fronteras sobre los ejes de cambio,
        núcleo protegido por la Regla de Dependencia y bounded context como criterio de
        granularidad.
  \item \textbf{P5, \texttt{Kittens} y \texttt{Rides} como \emph{plugins}:} la solución
        de Martin, el OCP aplicado a componentes, el patrón microkernel y los adaptadores.
  \item \textbf{P3 y P4:} el bounded context y la granularidad apoyan la división en
        contextos de P3; las clases base abstractas y las factorías ilustran el DIP de P4.
\end{itemize}


.
% Título: se omite «en Richards & Ford» del título oficial porque el Kitty
% Fuente de trabajo: notas.md (misma carpeta).
% Citas: páginas impresas; Richards & Ford (2020) no tiene paginación y se
% cita por capítulo/sección. Todas las claves están en core.bib.
% ============================================================

% Grupo local: \emergencystretch evita líneas desbordadas sin afectar a otros capítulos.
\begingroup\emergencystretch=3em \hbadness=2000
\section*{Notas de investigación: Kitty Problem, OCP y Plugin Architecture}

\section{Síntesis}

El \emph{Kitty Problem} es un ejemplo de \textcite[242--243]{martin2017}, no de Richards y
Ford. Parte de un agregador de taxis, presentado antes en el capítulo sobre el Principio
de Sustitución de Liskov \parencite[80]{martin2017}, construido como microservicios
divididos por función: \texttt{TaxiUI}, \texttt{TaxiFinder}, \texttt{TaxiSelector} y
\texttt{TaxiDispatcher} \parencite[242]{martin2017}. Cuando se pide un servicio de entrega
de gatitos, con restricciones por alergias de conductores y clientes, hay que modificar
\emph{todos} los servicios y coordinar su despliegue \parencite[243]{martin2017}. Martin
usa el caso para desmontar dos supuestos beneficios de los servicios: el desacoplamiento
y el desarrollo y despliegue independientes \parencite[240--241]{martin2017}. Su
diagnóstico es que las descomposiciones funcionales son muy vulnerables a los aspectos
transversales \parencite[243]{martin2017}.

La salida que propone no es crear más servicios, sino aplicar SOLID dentro de ellos
\parencite[244]{martin2017}. La lógica genérica queda en clases base abstractas; lo
específico de los viajes se extrae a un componente \texttt{Rides} y la nueva funcionalidad
va a un componente \texttt{Kittens}. Ambos extienden esas abstracciones mediante
\emph{Template Method} o \emph{Strategy}, respetan la Regla de Dependencia y son creados
por factorías controladas por la interfaz de usuario. Solo cambia \texttt{TaxiUI}; el
resto se extiende añadiendo un componente que se carga en tiempo de ejecución
\parencite[244]{martin2017}, lo que es el Principio Abierto/Cerrado aplicado a
componentes \parencite[246]{martin2017}. La lección arquitectónica: las fronteras no
pasan \emph{entre} los servicios, sino \emph{a través} de ellos
\parencite[246--247]{martin2017}.

Completan el cuadro tres ideas. El CCP es el SRP llevado a componentes, y Martin lo
vincula expresamente con el \emph{Kitty Problem} \parencite[105--106]{martin2017}. La
arquitectura de \emph{plugins} aísla las reglas de negocio de los componentes opcionales
o intercambiables \parencite[170--173]{martin2017}. Y
\textcite[caps.~8 y 17]{richardsford2020} aportan la partición por dominio, el antipatrón
\emph{entity trap} y el bounded context como criterio de granularidad.

\section{Conceptos clave}

\begin{table}[htbp]
  \centering
  \footnotesize
  \renewcommand{\arraystretch}{1.2}
  \begin{tabular}{@{}p{0.21\linewidth}p{0.47\linewidth}p{0.24\linewidth}@{}}
    \toprule
    \textbf{Concepto} & \textbf{Definición} & \textbf{Fuente} \\
    \midrule
    \raggedright \emph{Kitty Problem}
      & \raggedright Un requisito transversal obliga a cambiar y desplegar de forma
        coordinada todos los servicios de un sistema descompuesto por función
      & \raggedright \textcite[242--243]{martin2017} \tabularnewline
    \raggedright Falacias de los servicios
      & \raggedright Siguen acoplados por los datos que comparten; en esa medida, su
        desarrollo y despliegue deben coordinarse
      & \raggedright \textcite[240--241]{martin2017} \tabularnewline
    \raggedright OCP
      & \raggedright Abierto a la extensión y cerrado a la modificación (Meyer, citado por
        Martin); en arquitectura, con SRP y DIP
      & \raggedright \textcite[70, 73--74]{martin2017} \tabularnewline
    \raggedright CCP y cierre estratégico
      & \raggedright Reunir lo que cambia por las mismas razones; cerrar frente a los
        cambios más habituales
      & \raggedright \textcite[105--106]{martin2017} \tabularnewline
    \raggedright Arquitectura de \emph{plugins}
      & \raggedright Las reglas de negocio no dependen de los componentes opcionales; las
        fronteras siguen los ejes de cambio
      & \raggedright \textcite[170--173]{martin2017} \tabularnewline
    \raggedright Regla de Dependencia y adaptadores
      & \raggedright Las dependencias apuntan hacia adentro; los adaptadores convierten
        datos entre el núcleo y lo externo
      & \raggedright \textcite[203, 205]{martin2017} \tabularnewline
    \raggedright Servicios basados en componentes
      & \raggedright Estructura interna SOLID; lo nuevo entra como componentes nuevos y las
        fronteras atraviesan los servicios
      & \raggedright \textcite[245--247]{martin2017} \tabularnewline
    \raggedright Patrón microkernel
      & \raggedright Núcleo mínimo más \emph{plugins} independientes, localizados mediante
        un registro; adaptador para contratos propios
      & \raggedright \textcite[21--23]{richards2015} \tabularnewline
    \raggedright Partición por dominio y \emph{entity trap}
      & \raggedright Con partición técnica, un flujo se reparte entre capas; un componente
        por entidad confunde datos con flujos
      & \raggedright \textcite[cap.~8]{richardsford2020} \tabularnewline
    \raggedright Bounded context y granularidad
      & \raggedright Cada servicio modela un dominio o flujo; fronteras por propósito,
        transacciones y coreografía
      & \raggedright \textcite[caps.~7 y 17]{richardsford2020} \tabularnewline
    \bottomrule
  \end{tabular}
  \caption{Conceptos clave de la investigación INV-P5.}
  \label{tab:inv-p5-conceptos}
\end{table}

\section{Puntos de conexión con las preguntas}

\begin{itemize}
  \item \textbf{P5, por qué falló dividir por verbos o entidades:} la vulnerabilidad de la
        descomposición funcional, las falacias de los servicios, la partición técnica y el
        \emph{entity trap}, y la violación del CCP.
  \item \textbf{P5, diseño con DDD y Onion/Clean:} fronteras sobre los ejes de cambio,
        núcleo protegido por la Regla de Dependencia y bounded context como criterio de
        granularidad.
  \item \textbf{P5, \texttt{Kittens} y \texttt{Rides} como \emph{plugins}:} la solución
        de Martin, el OCP aplicado a componentes, el patrón microkernel y los adaptadores.
  \item \textbf{P3 y P4:} el bounded context y la granularidad apoyan la división en
        contextos de P3; las clases base abstractas y las factorías ilustran el DIP de P4.
\end{itemize}


